%%==========================================================================%%
%% validation.tex -- "Technical Validation".                                  %%
%% The section reviewers weight most. Demonstrates the data are FIT FOR        %%
%% PURPOSE, not a benchmark and not a results claim. Subsections:              %%
%% label-quality QC, inter-reader agreement, qualitative scoring, and a        %%
%% fit-for-purpose trained baseline. Point estimates are reported as        %%
%% fitness-for-purpose evidence attributed to the companion findings study. %%
%% The per-structure values and the confidence intervals are [TBD]. Report  %%
%% confidence intervals, never bare means. American spelling.               %%
%%==========================================================================%%

\section*{Technical Validation}\label{sec:validation}

The analyses in this section assess whether the released segmentation labels are of sufficient quality and consistency for the intended reuse, and whether the data support the development of a usable pediatric whole-body anatomy-segmentation model.
They are reported as evidence of fitness for purpose and not as scientific findings.

\subsection*{Label-quality and image-quality screening}\label{subsec:qc}
Each labeled examination was screened by an experienced reader to confirm that the imaging is of diagnostic quality and that the segmentation covers the intended structures without gross error.
Because the segmentations are produced and corrected by trained readers, this screen mainly excludes examinations whose image quality is too low to support reliable contouring rather than grading otherwise usable labels.
Examinations that failed the screen were excluded, and the counts are reported in the Participants subsection (Section~\ref{sec:methods}).
[TBD: report the number screened and the number excluded.]

\subsection*{Inter-reader agreement}\label{subsec:agreement}
To quantify the reliability of the manual labels, a subset of examinations was segmented independently by more than one reader, and agreement between readers was measured.
Per-structure overlap between readers is summarized with the Dice similarity coefficient, reported as the median and interquartile range with a case-level bootstrap 95\% confidence interval rather than as a pooled mean, so that the spread across examinations is retained.
Chance-corrected agreement is reported per structure as the Gwet AC2 coefficient with ordinal weights and a 95\% confidence interval.
Gwet AC2~\cite{Gwet_2008} is preferred over the Krippendorff alpha coefficient because the quality ratings are high in agreement and skewed in prevalence, the regime in which alpha and the kappa statistics collapse toward zero and understate concordance.
Figure~\ref{fig:agreement} shows the per-structure agreement, and Table~\ref{tab:metrics} reports the values.
[TBD: planned study, the number of readers and the number of multiply-annotated examinations, with the per-structure Dice and Gwet AC2 values and their confidence intervals.]

\subsection*{Qualitative rating}\label{subsec:likert}
In addition to the overlap analysis, two raters scored label quality on a five-point ordinal Likert scale, ranging from 1, unusable, to 5, clinically acceptable with no edits.
Ratings are reported per structure as the median and the category-count distribution across the five points.
Agreement between the two raters is reported per structure as the Gwet AC2 coefficient with quadratic weights and a 95\% confidence interval.
[TBD: planned, the per-structure rating distribution and the per-structure agreement values with confidence intervals.]

\subsection*{Fit-for-purpose segmentation baseline}\label{subsec:baseline}
To show that the released labels are sufficient to train a usable pediatric whole-body MRI anatomy-segmentation model, a baseline model was trained on the released labels with a self-configuring segmentation framework~\cite{Isensee_2021} and evaluated on a held-out portion of the subset (n = 5).
Across the held-out structures, the model trained on the released labels reached a mean Dice similarity coefficient of 78.2 and a median 95th-percentile Hausdorff distance (HD95) of 6.2 mm (interquartile range 4.4 to 10.0), an agreement level consistent with a usable pediatric whole-body anatomy-segmentation model.
Mean Dice is reported with a case-level bootstrap 95\% confidence interval, and HD95 is reported as the median and interquartile range with a bootstrap confidence interval on the median.
These figures originate in the separate study that developed the model on this dataset [TBD: companion findings publication, citation pending], and they are presented here only as evidence that the data are fit for the intended segmentation task, not as a benchmark.
With a held-out set of five examinations the intervals are wide, so the figures describe fitness for purpose rather than supporting an inferential benchmark, which belongs to that separate findings study.
The values are reported in Table~\ref{tab:metrics} and illustrated qualitatively in Figure~\ref{fig:baseline}.
[TBD: per-structure Dice and HD95 values with confidence intervals for the baseline, and the availability of the trained model and the training and evaluation scripts.]

%%-- Per-structure metric table placeholder -------------------------------%%
\begin{table}[h]
\caption{Label-quality and fit-for-purpose metrics per structure. Inter-reader agreement is the Dice similarity coefficient (median and interquartile range) and the Gwet AC2 coefficient between independent readers. Baseline performance is the Dice similarity coefficient and the 95th-percentile Hausdorff distance (HD95) of a model trained on the released labels. All values are reported with bootstrap 95\% confidence intervals [TBD: complete from the analysis].}\label{tab:metrics}
\begin{tabular}{@{}lcccc@{}}
\toprule
 & \multicolumn{2}{c}{Inter-reader agreement} & \multicolumn{2}{c}{Baseline (trained on released labels)} \\
\cmidrule(lr){2-3}\cmidrule(lr){4-5}
Structure & Dice & Gwet AC2 & Dice & HD95 [mm] \\
\midrule
Femur (L/R)           & [TBD] & [TBD] & [TBD] & [TBD] \\
Tibia (L/R)           & [TBD] & [TBD] & [TBD] & [TBD] \\
Fibula (L/R)          & [TBD] & [TBD] & [TBD] & [TBD] \\
Humerus (L/R)         & [TBD] & [TBD] & [TBD] & [TBD] \\
Pelvis (L/R)          & [TBD] & [TBD] & [TBD] & [TBD] \\
Spine                 & [TBD] & [TBD] & [TBD] & [TBD] \\
Sacrum                & [TBD] & [TBD] & [TBD] & [TBD] \\
Lung (L/R)            & [TBD] & [TBD] & [TBD] & [TBD] \\
Kidney (L/R)          & [TBD] & [TBD] & [TBD] & [TBD] \\
Liver                 & [TBD] & [TBD] & [TBD] & [TBD] \\
Spleen                & [TBD] & [TBD] & [TBD] & [TBD] \\
Heart                 & [TBD] & [TBD] & [TBD] & [TBD] \\
Urinary bladder       & [TBD] & [TBD] & [TBD] & [TBD] \\
\botrule
\end{tabular}
\end{table}

\subsection*{Metadata and data integrity}\label{subsec:integrity}
The integrity and internal consistency of the release were verified by schema validation of the metadata table, checks of referential consistency between the metadata and the imaging and label files, and per-file checksums.
Each released segmentation object was also checked for conformance to the DICOM-SEG standard and for referential integrity, confirming that every segmentation object correctly references its source image series.
Beyond these checks, the data pass the curation and validation that The Cancer Imaging Archive performs as part of submission.

%%-- Inter-reader agreement figure placeholder ----------------------------%%
\begin{figure}[t]
\centering
\fbox{\parbox[c][5cm][c]{0.85\textwidth}{\centering\itshape
[Figure 3 placeholder: per-structure inter-reader agreement, the Dice similarity
coefficient (median and interquartile range) and the Gwet AC2 coefficient, with
error bars showing bootstrap 95\% confidence intervals. Stat treatment to be
stated in the legend. Real figure pending the inter-reader study.]}}
\caption{\textbf{Inter-reader agreement of the manual anatomy segmentations.}
Per-structure agreement between independent readers, measured by the Dice similarity coefficient (median and interquartile range) and the Gwet AC2 coefficient, with error bars showing bootstrap 95\% confidence intervals. [TBD: state the number of readers, the number of multiply-annotated examinations, and the statistical treatment in the legend.]}
\label{fig:agreement}
\end{figure}

%%-- Baseline qualitative figure placeholder ------------------------------%%
\begin{figure}[t]
\centering
\fbox{\parbox[c][5cm][c]{0.85\textwidth}{\centering\itshape
[Figure 4 placeholder: qualitative segmentation surfaces on held-out
examinations produced by the model trained on the released labels, framed as
fitness-for-purpose evidence. Real figure pending the baseline predictions.]}}
\caption{\textbf{Qualitative segmentation surfaces on held-out examinations.}
Representative segmentation surfaces produced by a model trained on the released labels, illustrating that the released labels support a usable pediatric whole-body anatomy-segmentation model. [TBD: finalize panel content.]}
\label{fig:baseline}
\end{figure}
